% GENERATED FILE. Edit canonical lessons, then run npm run generate:books.
% canonical-source-hash: fnv1a64:e613f9b0a8458f1e
% canonical-lessons: RU-C95-bathtub, RU-C95-shelf, RU-C95-curtain, RU-C95-pillow, RU-C95-blanket

\chapter{Bathtub, Shelf, Curtain, Pillow, Blanket}
\label{ch:ru-bathtub-shelf-curtain-pillow-blanket}
\hlchaptermodality{95}

\begin{chapteropening}
\textbf{By the end of this chapter:} \emph{I can say a bathtub, a shelf, a curtain, a pillow and a blanket.}

Complete the last lesson of chapter 95: I can say a bathtub, a shelf, a curtain, a pillow and a blanket.
\end{chapteropening}

\section[\texorpdfstring{\ru{ванна} (vánna)}{ванна (vánna)}]{\ru{ванна} (vánna) — a bathtub}

\label{lesson:RU-C95-bathtub}



\begin{quote}
Before the new one: say the Russian for a cooker, then the Russian for a shower.
\end{quote}

\subsection*{You'll want to know: \ru{ванна}}
\textbf{\ru{ванна}} (\emph{vánna}) — \textquotedblleft{}a bathtub\textquotedblright{}.

\begin{cousinweb}
\textbf{\ru{ванна}} comes from German \emph{Wanne}, \textquotedblleft{}a tub\textquotedblright{}, or French \emph{vanne} (Vasmer); Trubachev notes it by 1697.
\end{cousinweb}

\begin{grammarlens}[title={What you've built}]
One more word for this chapter.
\end{grammarlens}

\subsection*{Your turn}
\begin{itemize}
\raggedright
  \item \emph{Say it:} \emph{vánna}
  \item \emph{Say it:} \emph{vánna}, once more
  \item \emph{Say it:} say \emph{vánna}
  \item \emph{Recall:} say the Russian for a cooker, then the Russian for a shower, then say \emph{vánna} again
\end{itemize}

\begin{tcolorbox}[breakable,colback=teal!4,colframe=teal!35!black,title={Before you move on}]
What does \ru{ванна} mean? (\textquotedblleft{}A bathtub\textquotedblright{}.) Say it once more.
\end{tcolorbox}
\section[\texorpdfstring{\ru{полка} (pólka)}{полка (pólka)}]{\ru{полка} (pólka) — a shelf}

\label{lesson:RU-C95-shelf}



\begin{quote}
Before the new one: say the Russian for a shower, then the Russian for a bathtub.
\end{quote}

\subsection*{You'll want to know: \ru{полка}}
\textbf{\ru{полка}} (\emph{pólka}) — \textquotedblleft{}a shelf\textquotedblright{}.

\begin{cousinweb}
\textbf{\ru{полка}} comes from \emph{\ru{полица}}, \textquotedblleft{}a shelf\textquotedblright{}, built on \emph{\ru{пол}}, \textquotedblleft{}floor, board\textquotedblright{} (Vasmer), which Vasmer relates to Sanskrit \emph{phálakam}, \textquotedblleft{}a board, a plank\textquotedblright{}.
\end{cousinweb}

\begin{grammarlens}[title={What you've built}]
One more word for this chapter.
\end{grammarlens}

\subsection*{Your turn}
\begin{itemize}
\raggedright
  \item \emph{Say it:} \emph{pólka}
  \item \emph{Say it:} \emph{pólka}, once more
  \item \emph{Say it:} say \emph{pólka}
  \item \emph{Recall:} say the Russian for a shower, then the Russian for a bathtub, then say \emph{pólka} again
\end{itemize}

\begin{tcolorbox}[breakable,colback=teal!4,colframe=teal!35!black,title={Before you move on}]
What does \ru{полка} mean? (\textquotedblleft{}A shelf\textquotedblright{}.) Say it once more.
\end{tcolorbox}
\section[\texorpdfstring{\ru{занавеска} (zanavéska)}{занавеска (zanavéska)}]{\ru{занавеска} (zanavéska) — a curtain}

\label{lesson:RU-C95-curtain}



\begin{quote}
Before the new one: say the Russian for a bathtub, then the Russian for a shelf.
\end{quote}

\subsection*{You'll want to know: \ru{занавеска}}
\textbf{\ru{занавеска}} (\emph{zanavéska}) — \textquotedblleft{}a curtain\textquotedblright{}.

\begin{cousinweb}
\textbf{\ru{занавеска}} is \emph{\ru{занавес}}, \textquotedblleft{}a curtain\textquotedblright{}, with \emph{-\ru{ка}}, \textquotedblleft{}little\textquotedblright{}. \emph{\ru{занавес}} is \emph{\ru{за}-}, \textquotedblleft{}across\textquotedblright{}, with \emph{\ru{навес}}, \textquotedblleft{}what is hung\textquotedblright{}, from \emph{\ru{вешать}}, \textquotedblleft{}to hang\textquotedblright{}.
\end{cousinweb}

\begin{grammarlens}[title={What you've built}]
One more word for this chapter.
\end{grammarlens}

\subsection*{Your turn}
\begin{itemize}
\raggedright
  \item \emph{Say it:} \emph{zanavéska}
  \item \emph{Say it:} \emph{zanavéska}, once more
  \item \emph{Say it:} say \emph{zanavéska}
  \item \emph{Recall:} say the Russian for a bathtub, then the Russian for a shelf, then say \emph{zanavéska} again
\end{itemize}

\begin{tcolorbox}[breakable,colback=teal!4,colframe=teal!35!black,title={Before you move on}]
What does \ru{занавеска} mean? (\textquotedblleft{}A curtain\textquotedblright{}.) Say it once more.
\end{tcolorbox}
\section[\texorpdfstring{\ru{подушка} (podúshka)}{подушка (podúshka)}]{\ru{подушка} (podúshka) — a pillow}

\label{lesson:RU-C95-pillow}



\begin{quote}
Before the new one: say the Russian for a shelf, then the Russian for a curtain.
\end{quote}

\subsection*{You'll want to know: \ru{подушка}}
\textbf{\ru{подушка}} (\emph{podúshka}) — \textquotedblleft{}a pillow\textquotedblright{}.

\begin{cousinweb}
\textbf{\ru{подушка}} $\leftarrow$ Proto-Slavic \emph{*podušĭka}. Vasmer reports the usual link with \emph{\ru{дух}}, \textquotedblleft{}breath\textquotedblright{}: a pillow as \textquotedblleft{}a thing blown up\textquotedblright{}. Machek's reading, \emph{\ru{под} \ru{ухо}}, \textquotedblleft{}under the ear\textquotedblright{}, has been rejected by others.
\end{cousinweb}

\begin{grammarlens}[title={What you've built}]
One more word for this chapter.
\end{grammarlens}

\subsection*{Your turn}
\begin{itemize}
\raggedright
  \item \emph{Say it:} \emph{podúshka}
  \item \emph{Say it:} \emph{podúshka}, once more
  \item \emph{Say it:} say \emph{podúshka}
  \item \emph{Recall:} say the Russian for a shelf, then the Russian for a curtain, then say \emph{podúshka} again
\end{itemize}

\begin{tcolorbox}[breakable,colback=teal!4,colframe=teal!35!black,title={Before you move on}]
What does \ru{подушка} mean? (\textquotedblleft{}A pillow\textquotedblright{}.) Say it once more.
\end{tcolorbox}
\section[\texorpdfstring{\ru{одеяло} (odeyálo)}{одеяло (odeyálo)}]{\ru{одеяло} (odeyálo) — a blanket}

\label{lesson:RU-C95-blanket}



\begin{quote}
Before the new one: say the Russian for a curtain, then the Russian for a pillow.
\end{quote}

\subsection*{You'll want to know: \ru{одеяло}}
\textbf{\ru{одеяло}} (\emph{odeyálo}) — \textquotedblleft{}a blanket\textquotedblright{}.

\begin{cousinweb}
\textbf{\ru{одеяло}} comes from Old Russian \emph{\ru{одеяти}}, \textquotedblleft{}to cover, to clothe\textquotedblright{}, the repeated-action form of \emph{\ru{одеть}}, \textquotedblleft{}to dress\textquotedblright{} (Vasmer): a blanket is what covers you.
\end{cousinweb}

\begin{grammarlens}[title={What you've built}]
One more word for this chapter.
\end{grammarlens}

\subsection*{Your turn}
\begin{itemize}
\raggedright
  \item \emph{Say it:} \emph{odeyálo}
  \item \emph{Say it:} \emph{odeyálo}, once more
  \item \emph{Say it:} say \emph{odeyálo}
  \item \emph{Recall:} say the Russian for a curtain, then the Russian for a pillow, then say \emph{odeyálo} again
\end{itemize}

\begin{tcolorbox}[breakable,colback=teal!4,colframe=teal!35!black,title={Before you move on}]
What does \ru{одеяло} mean? (\textquotedblleft{}A blanket\textquotedblright{}.) Say it once more.
\end{tcolorbox}
